\section{The Exponential Map}

\subsection{The exponential map}

\bdefn

    Let $G$ be a Lie group, a {\emphcolor one-parameter subgroup} is a Lie group morphism $\gamma\colon\bR\to G$.

\edefn

Note that a one-parameter subgroup is not a subgroup of $G$, it is a morphism, but its image is a subgroup (but we care about $\gamma$ too).

\bthrm

    Let $G$ be a Lie group, then the one-parameter subgroups of $G$ are precisely the maximal integral curves of left-invariant vector fields
    starting at $e\in G$.

\ethrm

\bproof

    If $\gamma$ is a maximal integral curve of $X\in\lie(G)$ starting at $e$, then since left-invariant vector fields are complete,
    $\gamma$ is defined on all of $\bR$.
    Since $X$ is $L_g$-related to itself for all $g\in G$, this means that $L_g\circ\gamma$ is integral for $X$ starting at $g$.
    In particular $L_{\gamma(s)}\circ\gamma$ is integral for $X$ starting at $\gamma(s)$.
    But $t\mapsto\gamma(t+s)$ is as well, so $\gamma(t+s)=\gamma(s)\circ\gamma(t)$, meaning $\gamma$ is a Lie group morphism.

    Conversely if $\gamma$ is a one-parameter subgroup, let $X=\gamma_*(d/dt)$ be the pushforward of $d/dt\in\X(\bR)$.
    Since $d/dt$ is left-invariant on $\bR$, $X$ is left-invariant on $G$.
    We just need to show that $\gamma$ is integral for $X$.
    Now,
    $$ \gamma'(t_0) = d\gamma_{t_0}\parens{\evdrv t_{t_0}} = X_{\gamma(t_0)} $$
    as desired.
    \qed

\eproof

Given $X\in\lie(G)$, the one-parameter subgroup of $G$ defined as above is said to be {\emphcolor generated} by it.
Because $X$ is determined by its value at the identity, a one-parameter subgroup is determined by its velocity at $e$.
So we have correspondences
$$ \set{\hbox{one-parameter subgroups of $G$}} \oto \lie(G) \oto T_eG . $$

\bprop

    For any $A\in\lgl_n(\bR)$ define
    $$ e^A = \sum_{k=0}^\infty\frac1{k!}A^k . $$
    This converges to an invertible matrix in $\gl_n(\bR)$, and the one-parameter subgroup of $\gl_n(\bR)$ generated by $A$
    is $\gamma(t)=e^{tA}$.

\eprop

\bproof

    See homework.
    \qed

\eproof

\bprop

    Let $H\subseteq G$ be a Lie subgroup.
    Then the one-parameter subgroups of $H$ are precisely the one-parameter subgroups of $G$ whose initial velocities lie in $T_eH$.

\eprop

\bproof

    If $\gamma\colon\bR\to H$ is a one-parameter subgroup, then composing with the inclusion $H\inj G$ gives a one-parameter subgroup of $G$
    whose initial velocity is $d\iota_e(\gamma'(0))\in d\iota_e(T_eH)$ which we identify with $T_eH\subseteq T_eG$.

    Conversely if $\gamma\colon\bR\to G$ is a one-parameter subgroup whose initial velocity lies in $T_eH$, let $\tilde\gamma\colon\bR\to H$
    be the unique one-parameter subgroup whose initial velocity is the same as $\gamma$'s.
    By the previous paragraph, composing $\iota\circ\tilde\gamma$ gives a one-parameter subgroup of $G$ with the same initial velocity as
    $\gamma$.
    By uniqueness, $\iota\circ\tilde\gamma=\gamma$.
    \qed

\eproof

\bexam

    For a Lie subgroup $H$ of $\gl_n(\bR)$, its one-parameter subgroups are curves of the form $e^{tA}$ where $A\in\lie(H)$ (since the initial
    velocity of $e^{tA}$ is $A$).

\eexam

\bdefn

    Let $G$ be a Lie group with Lie algebra $\g$.
    The {\emphcolor exponential map} of $G$ is the map $\exp\colon\g\to G$ defined by
    $$ \exp X = \gamma(1) $$
    where $X\in\g$ is left-invariant and $\gamma$ is the one-parameter subgroup generated by $X$.

\edefn

\bprop

    Let $G$ be a Lie group, then for any $X\in\lie(G)$, $\gamma(t)=\exp(tX)$ is the one-parameter subgroup of $G$ generated by $X$.

\eprop

\bproof

    Let $\gamma\colon\bR\to G$ be the one-parameter subgroup of $G$ generated by $X$, so by the rescaling lemma
    $\tilde\gamma\colon t\mapsto\gamma(ts)$ is the integral curve of $sX$ starting at $e$.
    Thus
    $$ \exp(sX) = \tilde\gamma(1) = \gamma(s) , $$
    as desired.
    \qed

\eproof

\bexam

    For $G=\gl_n(\bR)$ this shows that the exponential map is $\exp(A)=e^A$, explaining the name.

\eexam

\bprop

    Let $G$ be a Lie group with Lie algebra $\g$.
    \benum
        \item The exponential map is smooth ($\g$ is Euclidean).
        \item For any $X\in\g$ and $s,t\in\bR$, $\exp(s+t)X=\exp(sX)\exp(tX)$.
        \item For any $X\in\g$, $(\exp X)^{-1}=\exp(-X)$.
        \item For any $X\in\g$ and $n\in\bZ$, $(\exp X)^n=\exp(nX)$.
        \item The differential $(d\exp)_0\colon T_0\g\to T_eG$ is the identity map, under the canonical identification of $T_0\g$ and $T_eG$
        with $\g$.
        \item The exponential restricts to a diffeomorphism from a neighborhood of $0\in\g$ to a neighborhood of $e\in G$.
        \item If $H$ is another Lie group with Lie algebra $\h$, and $\Phi\colon G\to H$ a Lie group morphism, the diagram commutes

        \commsq{\g&\h\cr G&H}
            {\Phi_*}{\Phi}{\exp}{\exp}

        In other words, $\exp$ is a natural map from the Lie algebra functor (mapping $G\mapsto\g$) to the identity functor of Lie groups.
        \item The flow $\theta$ of $X\in\g$ is given by $\theta_t=R_{\exp(tX)}$ (right multiplication by $\exp(tX)$).
    \eenum

\eprop

\bproof

    For $X\in\g$ denote its flow by $\theta_{(X)}$.
    For the first point, we need to show that $\theta_{(X)}^{(e)}(1)$ depends smoothly on $X$.
    Define a vector field $\Xi$ on $G\times\g$ by
    $$ \Xi_{(g,X)} = (X_g,0) \in T_gG\oplus T_X\g \cong T_{(g,X)}(G\times\g) . $$
    This is smooth, as we show now.
    Let $X_1,\dots,X_n$ be a basis for $\g$ and let $(x^i)$ be the corresponding global coordinates of $\g$ (the coordinates represented by this
    basis).
    Let $(w^i)$ be any local coordinates for $G$, then for $f\in C^\infty(G\times\g)$ we have
    $$ \Xi f(w^i,x^i) = \Xi_{(w^i,X^i)}f = ((x^iX_i)_{w^i},0)f = x^iX_if(w^j) . $$
    This is clearly a smooth function, so $\Xi$ is smooth.

    Let $\Theta$ be $\Xi$'s flow, we can see that
    $$ \Theta_t(g,X) = (\theta_{(X)}(t,g),X) , $$
    since differentiating this with respect to $t$ gives $((\theta_{(X)}^{(g)})'(t_0),0)=(X_{\theta_{(X)}^{(g)}(t_0)},0)$ which is equal
    to $\Xi$ at $\Theta_{t_0}(g,X)$.
    This is smooth as flow is smooth, and so $\Theta_1(e,X)=(\theta_{(X)}^{(e)}(1),X)$ is smooth, as desired.

    The next three points are clear, since $t\mapsto\exp(tX)$ is a one-parameter subgroup, so a group morphism.

    To prove (5), let $X\in\g$ be arbitrary and $\sigma\colon\bR\to\g$ be given by $\sigma(t)=tX$.
    Then $\sigma'(0)=X$ (under the identification $T_0\g\cong\g$, $\sigma$ and $X$ are identified),
    $$ (d\exp)_0(X) = (d\exp)_0(\sigma'(0)) = (\exp\circ\sigma)'(0) = \evdrv t_{t=0}\exp(tX) = X , $$
    since $\exp(tX)$ is the one-parameter subgroup of $X$.

    Part (6) is due to the inverse function theorem.

    To prove naturality we need to show $\exp(\Phi_*X)=\Phi(\exp X)$ for all $X\in\g$.
    We show the more general
    $$ \exp(t\Phi_*X) = \Phi(\exp(tX)) $$
    for all $t\in\bR$.
    The left-hand side is the one-parameter subgroup generated by $\Phi_*X$.
    Let $\sigma(t)=\Phi(\exp(tX))$; we want to show that it is also the one-parameter subgroup of $\Phi_*X$, i.e.\ it is a Lie group
    morphism satisfying $\sigma'(0)=(\Phi_*X)_e$.
    As the composition of Lie group morphisms, it itself is a Lie group morphism.
    And
    $$ \sigma'(0) = d\Phi_e\parens{\evdrv t_{t=0}(\exp(tX)} = d\Phi_e(X_e) = (\Phi_*X)_e . $$

    For the final point we want to show $(\theta_{(X)})_t=R_{\exp(tX)}$.
    We use the fact that $L_g$ maps integral curves of $X$ to integral curves of $X$, so $t\mapsto L_g(\exp(tX))$ is an integral curve of $X$
    starting at $g$, so it is equal to $\theta_{(X)}^{(g)}$.
    Thus
    $$ R_{\exp(tX)}(g) = g\exp(tX) = L_g(\exp(tX)) = \theta^{(g)}_{(X)}(t) = (\theta_{(X)})_t(g) $$
    as desired.
    \qed

\eproof

\bnote

    Importantly, the above proposition does not imply $\exp(X+Y)=\exp(X)\exp(Y)$.
    As it turns out, when $G$ is connected this is only true if it is Abelian.

\enote

\subsection{Normal subgroups}

\blemm

    Let $G$ be a connected Lie group, and $H\subseteq G$ a connected Lie subgroup.
    Let $\g,\h$ be their respective Lie algebras, then $H$ is normal in $G$ iff for all $X\in\g$ and $Y\in\h$,
    $$ \exp(X)\exp(Y)\exp(-X) \in H . $$

\elemm

\bproof

    Since $\exp(-X)=\exp(X)^{-1}$, if $H$ is normal this holds.
    Conversely, let $V\subseteq\g$ be a neighborhood of $0$ and $W\subseteq G$ such that $\exp\colon V\to W$ is a diffeomorphism.
    Since the exponential of $H$ is the restriction of that of $G$, shrinking $V$ if necessary, the restriction of $\exp$ to $V\cap\h$
    is a diffeomorphism to a neighborhood $U_0\subseteq H$ of $e$.
    We can further assume that $V$ is symmetric (if $X\in V$ then $-X\in V$) by potentially further shrinking $V$.

    This implies that $ghg^{-1}\in H$ whenever $g\in U$ and $h\in U_0$ (since they are in the image of $\exp$).
    Recall that we showed that $U_0$ and $U$ generate $H,G$ respectively we get the desired result.
    \qed

\eproof

\bdefn

    If $\g$ is a Lie algebra, a {\emphcolor kernel} is a subspace $\h\subseteq\g$ which is closed under Lie bracketing from $\g$:
    for $X\in\g$ and $Y\in\h$, $[X,Y]\in\h$.

\edefn

For a Lie group $G$ and $g\in G$, the map $C_g\colon G\to G$ given by conjugation with $g$: $C_g(h)=ghg^{-1}$ is a Lie group
isomorphism, and so if we define its induced map $\Ad(g)=(C_g)_*\colon\g\to\g$.
This defines a map
$$ \Ad\colon G\to\gl(\g) $$
called the {\emphcolor Adjoint representation} of $G$.

\bprop

    The adjoint representation is indeed a representation.

\eprop

\bproof

    It is well-defined since $C_g$ are Lie group isomorphisms so their induced maps are linear automorphisms of $\g$.
    Since $C_{g_1g_2}=C_{g_1}\circ C_{g_2}$ we see that $\Ad(g_1g_2)=\Ad(g_1)\circ\Ad(g_2)$, so it is a morphism.
    We need to show only that $\Ad$ is smooth.

    Let $C\colon G\times G\to G$ be given by $C(g,h)=ghg^{-1}$, which is smooth.
    Let $X\in\g$ and $g\in G$, so $\Ad(g)X$ is the left-invariant vector field whose value at $e\in G$ is
    $$ (\Ad(g)X)_e = ((C_g)_*X)_e = \evdrv t_{t=0}C_g(\exp tX) = \evdrv t_{t=0}C(g,\exp tX) = dC_{(g,e)}(0,X_e) . $$
    Because $dC$ is smooth, this is smoothly dependent on $X$ and $g$.
    If we choose a basis $(E_i)$ for $\g$, then $\gl(\g)$ has a basis with respect to this.
    Let $(\epsilon^j)$ be the dual basis of $(E_i)$, then the matrix entries of $\Ad(g)$ is $(\Ad(g))^j_i=\epsilon^j(\Ad(g)E_i)$.
    Plugging in $X=E_i$ to the above we see that these coordinates are smooth.
    \qed

\eproof

\bdefn

    If $\g$ is a Lie algebra, define its {\emphcolor adjoint representation} $\ad\colon\g\to\lgl(\g)$ by
    $$ \ad(X)Y = [X,Y] . $$
    Note that its codomain is $\lgl(\g)$, not $\gl(\g)$.

\edefn

This is a smooth map, as it is a linear map between Euclidean spaces.

\bthrm

    Let $G$ be a Lie group with Lie algebra $\g$.
    Let $\Ad\colon G\to\gl(\g)$ be its adjoint representation, then its induced Lie algebra representation $\Ad_*\colon\g\to\lgl(\g)$
    is just $\Ad_*=\ad$.

\ethrm

\bproof

    Let $X\in\g$, then $\Ad_*X$ is determined by its value at the identity, an element of $\lgl(\g)$.
    Now, we have
    $$ (\Ad_*X)Y = \parens{\evdrv t_{t=0}\Ad(\exp tX)}Y = \evdrv t_{t=0}(\Ad(\exp tX)Y) . $$
    Now $C_g=R_{g^{-1}}\circ L_g$ so that $\Ad(g)=(R_{g^{-1}})_*\circ(L_g)_*$, and thus
    $$ (\Ad(\exp(tX))Y)_e = dR_{\exp(-tX)}\circ dL_{\exp(tX)}(Y_e) = dR_{\exp(-tX)}(Y_{\exp(tX)}) , $$
    since $Y$ is left-invariant.
    Recall that the flow of $X$ is given by $\theta_t(g)=R_{\exp(tX)}(g)$ so this is equal to
    $$ = d(\theta_{-t})(Y_{\theta_t(e)}) . $$
    Thus if we take the derivative of this with respect to $t$ and setting $t=0$, we get
    $$ ((\Ad_*X)Y)_e = \evdrv t_{t=0}d(\theta_{-t})(Y_{\theta_t(e)}) = (\L_XY)_e = [X,Y]_e . $$
    Since these vector fields are determined by their values at the identity, we see that $(\Ad_*X)Y=\ad(X)Y$ as desired.
    \qed

\eproof

\bthrm

    Let $H\subseteq G$ be a connected Lie subgroup, and denote the respective Lie algebras by $\h,\g$.
    Then $H$ is a normal subgroup of $G$ iff $\h$ is an ideal of $\g$.

\ethrm

\bproof

    Consider the commutative diagram

    \commsq{\g&\g\cr G&G}
        {\Ad(g)}{C_g}{\exp}{\exp}

    for all $g\in G$.
    If $\h$ is an ideal then for $Y\in\h$ and $g=\exp X$ we see that
    $$ \exp(\Ad(\exp X)Y) = C_{\exp X}(\exp Y) = (\exp X)(\exp Y)(\exp(-X)) . $$
    Conversely a similar commutative diagram for $\Ad_*=\ad$ yields
    $$ \Ad(\exp X) = \exp(\ad X) $$
    so that
    $$ \Ad(\exp X)Y = \exp(\ad X)Y = \sum_{k=0}^\infty\frac1{k!}(\ad X)^kY . $$
    Since $\h$ is an ideal, $(\ad X)Y=[X,Y]\in\h$, and by induction $(\ad X)^kY\in\h$.
    Furthermore $\h$ is closed in $\g$ so that $\Ad(\exp X)Y\in\h$, and thus its exponent is in $\h$.
    So we have $(\exp X)(\exp Y)(\exp(-X))\in\h$, so $H$ is normal.

    Conversely if $H$ is normal, then for $X\in\g$ and $Y\in\h$ then by the commutative diagram
    $$ \exp(\Ad(\exp(tX))sY) = (\exp tX)(\exp sY)(\exp tX)^{-1} \in H . $$
    Since $\Ad(\exp(tX))$ is linear over $\bR$, it follows that
    $$ \exp(s\Ad(\exp(tX))Y) = \exp(\Ad(\exp(tX))sY) , $$
    which is in $H$ for all $s$, and thus $\Ad(\exp(tX))Y\in\h$ (since its integral curve lies in $H$).
    By above,
    $$ \evdrv t_{t=0}\Ad(\exp(tX)Y) = [X,Y], $$
    so that $[X,Y]\in\h$ and thus $\h$ is an ideal.
    \qed

\eproof

\subsection{Bi-invariant metrics}

In this section we discuss metrics on Lie groups.
Since Riemannian metrics are typically denoted by $g$, we will in this section denote elements of a Lie group by $\phi,\psi$, etc.

\bdefn

    Let $G$ be a Lie group.
    A Riemmanian metric $g$ on $G$ is {\emphcolor left-invariant} if $L_\phi^*g=g$ for all $\phi\in G$.
    Similarly we define {\emphcolor right-invariant} Riemannian metrics, and {\emphcolor bi-invariant metrics} are metrics which are both
    left- and right-invariant.

\edefn

\blemm

    Let $G$ be a Lie group with Lie algebra $\g$.
    Then a Riemannian metric $g$ on $G$ is left-invariant iff for all $X,Y\in\g$ the function $g(X,Y)\in C^\infty(G)$ is constant.

\elemm

\bproof

    Recall that $g(X,Y)(\phi)=g_\phi(X_\phi,Y_\phi)$; we claim that if $g$ is left-invariant this is equal to $g_e(X_e,Y_e)$.
    Indeed, by left-invariance of $g,X,Y$,
    $$ g_e(X_e,Y_e) = (L_\phi^*g)_e(X_e,Y_e) = g_\phi(d(L_\phi)_eX_e,d(L_\phi)_eY_e) = g_\phi(X_\phi,Y_\phi), $$
    as desired.
    The same computation shows that $L_\phi^*g=g$ when $g(X,Y)$ is constant.
    \qed

\eproof

From this lemma we deduce that every left-invariant Riemannian metric arises uniquely from an inner product on $\g$.
That is, $g\mapsto g_e$ is a bijection.

\bdefn

    Let $V$ be a Euclidean space, $\gl(V)$ be its group of linear automorphisms.
    Let $H\subseteq\gl(V)$ be a subgroup, then we say an inner product $\iprod{\bul,\bul}$ is {\emphcolor $H$-invariant} if
    $\iprod{h\bul,h\bul}=\iprod{\bul,\bul}$ for all $h\in H$.

\edefn

\bprop

    Let $G$ be a Lie group with Lie algebra $\g$, and $g$ a left-invariant Riemannian metric on $G$.
    Let $\iprod{\bul,\bul}$ denote its associated inner product on $\g$ (i.e.\ $g_e$).
    Then $g$ is bi-invariant iff $\iprod{\bul,\bul}$ is $\Ad(G)\subseteq\gl(\g)$-invariant.

\eprop

\bproof

    For $\phi\in G$ recall that $C_\phi=R_{\phi^{-1}}\circ L_\phi$.
    So for $X\in\g$ we have that
    $$ \Ad(\phi)X = (R_{\phi^{-1}})_*(L_\phi)_*X = (R_{\phi^{-1}})_*X , $$
    by left-invariance of $X$.
    Thus for all $\psi\in G$ and $X,Y\in\g$,
    $$ (R_{\phi^{-1}}^*g)_\psi(X_\psi,Y_\psi) = g_{\psi\phi^{-1}}((R_{\phi^{-1}})_*X_{\psi\phi^{-1}},(R_{\phi^{-1}})_*Y_{\psi\phi^{-1}})
    = g_{\psi\phi^{-1}}(\Ad(\phi)X_{\psi\phi^{-1}},\Ad(\phi)Y_{\phi\psi^{-1}}) . $$
    By left-invariance of $g$ this is just equal to
    $$ \iprod{\Ad(\phi)X,\Ad(\phi)Y} . $$

    So if the inner product is $\Ad(G)$-invariant we see that $(R_{\phi^{-1}})^*g=g$ for all $\phi$ so that $g$ is right-invariant as
    well, so bi-invariant.

    Conversely if $g$ is bi-invariant then we have that $(R_{\phi^{-1}})^*g=g$.
    So setting $\psi=e$ we see that $g_e=\iprod{\Ad(\phi)\bul,\Ad(\phi)\bul}$ so that $g_e$ is $\Ad(G)$-invariant.
    \qed

\eproof

\blemm

    Let $V$ be a Euclidean space, and $H$ a subgroup of $\gl(V)$.
    Then there is an $H$-invariant inner product on $V$ iff $H$ has compact closed in $\gl(V)$.

\elemm

\bproof

    Omitted.
    \qed

\eproof

The following is then an immediate result of this lemma and the above proposition.

\bthrm

    A Lie group $G$ has a bi-invariant Riemannian metric iff $\Ad(G)$ has compact closure in $\gl(\g)$.

\ethrm

\bcoro

    Every compact Lie group has a bi-invariant Riemannian metric.

\ecoro

\bproof

    $\Ad(G)$ is the continuous image of $G$ under $\Ad$, and thus compact.
    \qed

\eproof

\bexam

    If $G$ is Abelian, then $C_g=\id$ and thus $\Ad$ is trivial and so every left-invariant metric is bi-invariant.

\eexam

